\subsection{Décomposition géométrique sous inversion de feuillet}
\label{sec:ii09-descomposicion-paridad}

Dans la coordinatisation angulaire, l’inversion de feuillet conserve \(a,d\) et
inverse \(u\). Le vecteur exprimé admet la décomposition
\[
 \boldsymbol\theta_+=a c_A+u c_h+d c_\Delta,
 \qquad
 \boldsymbol\theta_-=a c_A-u c_h+d c_\Delta.
\]
\begin{proposition}[Composantes paire et impaire du transport angulaire]
\label{prop:ii09-paridad-angular}
Les deux vecteurs satisfont
\[
 \frac{\boldsymbol\theta_++\boldsymbol\theta_-}{2}
       =a c_A+d c_\Delta,\qquad
 \frac{\boldsymbol\theta_+-\boldsymbol\theta_-}{2}=u c_h.
\]
La contribution qui change de signe occupe la droite
\(\operatorname{span}(c_h)\) ; la composante conservée occupe
\(\operatorname{span}(c_A,c_\Delta)\).
\end{proposition}
\begin{proof}
La somme annule les termes \(u c_h\), et la différence annule
les deux autres. Les colonnes sont linéairement indépendantes, puisque
\(\det\mathsf M_{\rm CKM}=-3857/6\). Les deux
composantes forment donc une décomposition directe.
\end{proof}

La même inversion agit sur l’évaluation de Barbero--Immirzi :
elle échange \(q_+\) et \(q_-\), et inverse la composante bilinéaire
dans \(A\) et \(C^*\). Par conséquent, la fonctionnelle d’incidence
est impaire. Une observable symétrique dans les deux canaux conserve sa
valeur. Cette différence de parité exprime comment une même
opération sur la structure produit des réponses différentes selon
le lecteur : orientation aréale, information symétrique et vecteur de
mélange conservent des informations différentes du même antécédent.
